% REPLACES: main.tex lines 1815-1880 ("\subsubsection{Logistic noise}" to "posterior mass sits relative to the common bounds \(m,M\).")
% RETAINS LABELS: logistic-noise, result-proposition-a-5, eq:paper-a-21,
% eq:paper-a-22, eq:paper-a-23. No label is removed.
% NEW SYMBOLS: mathcal C_L, L, mu_f(x;pi), bar mu_f(pi), mathcal E_f, mathsf E_f.
% The symbols denote the density class, score bound, full-order posterior,
% posterior supremum, favorable flow set, and entry. Existing alpha_j refers to this set.
% Uses Gamma and u_0 from the nonemptiness argument; both are defined again below.
% DEPENDS ON: the quality-class Proposition 2, prior pi, zero preparation cost,
% kappa(pi), and the general nonemptiness argument in the main-model revision.
% OWNERSHIP: use the exact condition below; the Laplace-specific sufficient
% condition based on alpha_H > 1/2 does not extend to every class member.
% NUMBERS: none. CITES: ref-Blackwell1953 already exists.
% SCOPE: the common-variance comparison below is a numerical diagnostic.

\subsubsection{Bounded-score noise}\label{logistic-noise}

Let \(\mathcal C_L\) denote the positive, absolutely continuous densities
\(f\) with \(|f'|\le Lf\) almost everywhere, where \(0<L<1\).
Equivalently, \(\log f\) is Lipschitz with constant \(L\).
Set \(b=1/L\).
The bounds \(m(\pi),M(\pi)\) and
\(\kappa(\pi)=\min\{m(\pi),1-M(\pi)\}\) use this value of \(b\).
For full orders, define

\[
\mu_f(x;\pi)=
\frac{\pi f(x-1)}{\pi f(x-1)+(1-\pi)f(x+1)},\qquad
\bar\mu_f(\pi)=\operatorname*{ess\,sup}_{x\in\mathbb R}\mu_f(x;\pi).
\]

The essential supremum uses Lebesgue measure.
The unconditional flow law and both conditional flow laws are equivalent to that measure.
The posterior is continuous, so this essential supremum equals \(\sup_{x\in\mathbb R}\mu_f(x;\pi)\).
The supremum satisfies \(\pi<\bar\mu_f(\pi)\le M(\pi)\).
It need not be attained at a finite flow.
It can lie strictly below \(M(\pi)\).

\Needspace{8\baselineskip}

\protect\phantomsection\label{result-proposition-a-5}{\textbf{Proposition
A.5 (bounded-score noise; analytical).}} \emph{Replace Laplace noise by
any \(f\in\mathcal C_L\).
Keep the quality-class assumptions of
\hyperref[result-proposition-2]{Proposition 2}, with preparation costs
\(C\in\{0,c_H\}\).
Replace the upper posterior in \eqref{eq:paper-a2} by \(\bar\mu_f(\pi)\), so that}

\[
B_{r_0}(\pi)<c_H<B_{r_1}(\bar\mu_f(\pi)).
\]

\emph{Retain \eqref{eq:paper-a3} with \(b=1/L\), namely}

\[
\Delta_T(r_0)<k<(1-L)\rho\kappa(\pi)\Delta_T(r_1).
\]

\emph{Then the weak economy has unique zero orders and entry \(\rho\).
The strong economy has unique full orders \((1,-1)\) and an informative price.
Let}

\[
\begin{aligned}
\tau&=\frac{c_H-g_L(r_1)}{g_H(r_1)-g_L(r_1)},
&\mathcal E_f&=\{x:\mu_f(x;\pi)\ge\tau\},\\
\alpha_H&=\int_{\mathcal E_f}f(x-1)\,dx,
&\alpha_L&=\int_{\mathcal E_f}f(x+1)\,dx.
\end{aligned}
\]

\emph{Total entry satisfies}

\[
\mathsf E_f=\rho+(1-\rho)[\pi\alpha_H+(1-\pi)\alpha_L]>\rho.
\]

\emph{High-class ownership rises exactly when
\([\rho+(1-\rho)\alpha_H]W_H(r_1)>\rho W_H(r_0)\).
The condition holds automatically in the point-mass case with \(W_H=1\).
Part (iii) of \hyperref[result-proposition-2]{Proposition 2} also holds at
\(r_2>r_1\) if \(B_{r_2}(M(\pi))<c_H\): orders remain \((1,-1)\), and entry equals \(\rho\).
For fixed \(f,L,\pi,\rho,p,G_H,G_L\) with
\(\mathbb E_H[\max\{\theta,p\}]>\mathbb E_L[\max\{\theta,p\}]\),
these comparisons and the ownership increase hold on a nonempty open set
of uniform strengths and costs \((r_0,r_1,r_2,c_H,k)\).
The other preparation cost stays zero.}

\emph{Proof.}
The score bound gives
\(|\log f(y)-\log f(z)|\le L|y-z|\).
Hence \(f(y)\ge f(z)e^{-L|y-z|}\).
Integration over \(y\) gives \(f(z)\le L/2\).
Also \(\int|f'|\le L\), so \(f\in W^{1,1}\).
The density is positive, bounded, and continuous.
The null-set arguments in Online Appendix A.1 apply under every mixed order and every unilateral deviation.
\par
For any feasible \(q,q'\), the density ratio lies in \([e^{-2L},e^{2L}]\).
Integration over the order distributions and Bayes' rule give
\(m(\pi)\le\mu_X\le M(\pi)\).
Conditional expectation gives the same bounds for the price posterior.

The zero preparation cost gives \(e\ge\rho\) at every belief.
The price inversion and the residual calculation in Online Appendix A.3 apply without change.
They use only \(t_0,t_j,g_j,\Delta_T\), independence, and the positive entry floor.
In particular,

\[
\rho\kappa(\pi)\Delta_T\le A_H,A_L\le\Delta_T.
\]

The translation identity, Fubini argument, and product rule in Online Appendix A.3 allow this general \(f\).
They give \(|F_j'|\le LF_j\) and

\[
U_j'(s)\ge(1-sL)F_j(s)-k
\ge(1-L)\rho\kappa(\pi)\Delta_T(r_1)-k>0.
\]

These steps differentiate the density.
They do not differentiate entry.
At \(r_0\), every correctly signed nonzero order earns at most
\(s[\Delta_T(r_0)-k]<0\).
A wrongly signed order loses as well.
Thus zero is the unique best response, including among mixed orders.
The posterior equals \(\pi\), and the lower cost restriction excludes expensive preparation.
The constant competitive price constructs the weak outcome.
\par
At \(r_1\), the positive derivative makes unit magnitude strictly optimal for each class.
It excludes every mixture with another order.
Define entry and price by \eqref{eq:paper-9}, with posterior \(\mu_f(x;\pi)\).
The price increases strictly in the posterior.
Its measurable inverse gives price sufficiency.
No monotonicity of the posterior in flow is required.
The construction gives competitive prices and optimal preparation.
The global bound verifies every order deviation.

Gross preparation profit falls with incumbent strength at every fixed belief.
The cost window therefore gives \(\pi<\tau<\bar\mu_f(\pi)\).
By the definition of the essential supremum, \(\{\mu_f>\tau\}\) has positive Lebesgue measure.
Both conditional densities are positive, so \(\alpha_H,\alpha_L>0\).
On \(\mathcal E_f\), Bayes' rule gives

\[
f(x-1)\ge
\frac{(1-\pi)\tau}{\pi(1-\tau)}f(x+1),\qquad
\alpha_H\ge\frac{(1-\pi)\tau}{\pi(1-\tau)}\alpha_L>\alpha_L.
\]

The entry formula and the stated ownership condition follow.
The likelihood ratio is not constant.
Both conditional densities integrate to one, so a constant ratio would equal one.
Then \(f\) would be two-periodic.
No positive integrable density is two-periodic.
The strong price experiment is informative.
It strictly Blackwell-dominates the weak constant experiment.
Both conditional densities integrate to one, so their nonconstant ratio exceeds one on a set of positive measure.
Hence \(\bar\mu_f(\pi)>\pi\).
\par
For part (iii) of \hyperref[result-proposition-2]{Proposition 2},
\(\Delta_T(r_2)\ge\Delta_T(r_1)\) preserves the strict order bound.
It forces full orders, and \(\mu_X\le M(\pi)\) excludes expensive preparation.

For nonemptiness, fix \(f,L,\pi,\rho,p,G_H,G_L\) under the stated assumptions.
As in \hyperref[nonemptiness-for-every-value-gap]{the nonemptiness argument}, put

\[
\Gamma=\int_p^\infty[G_L(u)-G_H(u)]\,du>0,\qquad
u_0=\inf\left\{x\ge p:\int_p^x[G_L(u)-G_H(u)]\,du>0\right\}.
\]

For uniform incumbents, that argument gives \(\Delta_T(u_0)=0\) and \(g_H(u_0)-g_L(u_0)=\Gamma\).
The spread is continuous and strictly increasing above \(u_0\).
For finite \(r\ge u_0\), \(g_H(r)-g_L(r)\ge p\Gamma/r>0\).
The functions \(g_j(r)\) and \(W_H(r)\) are continuous and nonincreasing, with \(W_H(u_0)>0\).

\begin{enumerate}
\item
Choose \(\tau\in(\pi,\bar\mu_f(\pi))\) at a continuity point of the posterior law in each class.
Each law has at most countably many atoms, so such a point exists.
This fixes \(\mathcal E_f,\alpha_H,\alpha_L\) before any strength or cost is chosen.
In particular, \(\alpha_H,\alpha_L>0\).

\item
Choose \(r_1>u_0\) sufficiently near \(u_0\) that
\[
B_{r_1}(\tau)>B_{u_0}(\pi),\qquad
[\rho+(1-\rho)\alpha_H]W_H(r_1)>\rho W_H(u_0).
\]
The first inequality follows from
\(B_{u_0}(\tau)-B_{u_0}(\pi)=(\tau-\pi)\Gamma>0\).
The second follows from \(\alpha_H>0\) and \(W_H(u_0)>0\).
Continuity permits the same \(r_1\) in both inequalities.

\item
Set \(c_H=B_{r_1}(\tau)=g_L(r_1)+\tau[g_H(r_1)-g_L(r_1)]\).
Thus the preparation threshold equals the \(\tau\) chosen in the first step.
The positive profit gap gives \(c_H<B_{r_1}(\bar\mu_f(\pi))\).

\item
Choose \(r_0\in(u_0,r_1)\) sufficiently near \(u_0\) that
\[
\Delta_T(r_0)<(1-L)\rho\kappa(\pi)\Delta_T(r_1).
\]
This is possible because \(\Delta_T(r_0)\to0\) as \(r_0\downarrow u_0\).
Monotonicity gives \(B_{r_0}(\pi)\le B_{u_0}(\pi)<c_H\).
It also gives \(W_H(r_0)\le W_H(u_0)\), so the strict ownership inequality remains valid.

\item
Choose \(k\) strictly between \(\Delta_T(r_0)\) and
\((1-L)\rho\kappa(\pi)\Delta_T(r_1)\).
Both the preparation-cost window and the order-cost window now hold for this single choice of primitives.

\item
Choose \(r_2>r_1\) sufficiently large that \(B_{r_2}(M(\pi))<c_H\).
Finite class means give \(g_j(r)\to0\), so such a strength exists.
The larger spread preserves the strict order bound at \(r_2\).
\end{enumerate}

All inequalities have strict slack at the constructed point.
The threshold is continuous in \((r_1,c_H)\), since the class profit gap is positive.
At the chosen threshold, both conditional entry probabilities are continuous in that threshold.
The strict ownership inequality therefore persists on a neighborhood of \((r_0,r_1,r_2,c_H,k)\).
Continuity preserves the other strict inequalities there as well.
This proves openness in those five coordinates, with \(f,L,\pi,\rho,p,G_H,G_L\) and zero preparation cost fixed.
\(\square\)

\emph{Logistic example.}
The density

\[
f_{\mathrm{log}}(z)=\frac{1}{4b}\operatorname{sech}^{2}\left(\frac z{2b}\right)
=\frac{e^{-z/b}}{b(1+e^{-z/b})^2},\qquad b>1,
\tag{A.21}\label{eq:paper-a-21}
\]

has log derivative \(-\tanh(z/2b)/b\), so it belongs to \(\mathcal C_{1/b}\).
Under full orders its posterior satisfies

\[
\operatorname{logit}\mu_X
=\operatorname{logit}\pi
+2\log\cosh\left(\frac{x+1}{2b}\right)
-2\log\cosh\left(\frac{x-1}{2b}\right).
\tag{A.22}\label{eq:paper-a-22}
\]

The derivative is
\([\tanh((x+1)/2b)-\tanh((x-1)/2b)]/b>0\).
The posterior spans \((m(\pi),M(\pi))\) and attains neither endpoint.
Thus \(\bar\mu_f(\pi)=M(\pi)\).
The original upper restriction applies.
\par
Set \(A=e^{1/b}\) and
\(w=\sqrt{\tau(1-\pi)/[(1-\tau)\pi]}\).
With \(y=e^{x/b}\), the square root of the likelihood ratio is
\((Ay+1)/(y+A)\).
The equation \(w=(Ay+1)/(y+A)\) gives

\[
\begin{aligned}
x^*_{\mathrm{log}}&=b\log\frac{Aw-1}{A-w},\\
\alpha_H^{\mathrm{log}}&=\frac1{1+e^{(x^*_{\mathrm{log}}-1)/b}},\qquad
\alpha_L^{\mathrm{log}}=\frac1{1+e^{(x^*_{\mathrm{log}}+1)/b}}.
\end{aligned}
\tag{A.23}\label{eq:paper-a-23}
\]

These formulas hold for \(m(\pi)<\tau<M(\pi)\).
At \(\pi=1/2\), they reduce to the benchmark formulas.
The class also contains hyperbolic-secant noise, suitably scaled Student noise, and Gaussian noise plus an independent Laplace component.
Take these examples symmetric about zero, as are the Laplace and logistic laws above.
Each then satisfies the median-zero condition in Online Appendix A.6.

Laplace and logistic noise are compared at a common scale parameter \(b\), not at a common variance.
At a common scale, the logistic law equals the Laplace law plus independent noise.
The Laplace flow experiment therefore Blackwell-dominates the logistic experiment at every order profile, pure or mixed
(\citeproc{ref-Blackwell1953}{Blackwell 1953}).
Under the price construction above, each price is a measurable function of flow and reveals its flow posterior.
Each price experiment is therefore Blackwell-equivalent to its flow experiment.
The same dominance holds for prices, including the full-order equilibria compared here.
Online Appendix A.6 proves the decomposition.
\par
At a common variance, neither experiment dominates in the numerical comparison.
That comparison is a numerical diagnostic.
Dominance does not order entry: \(\Pr(\mu_X\ge\tau)\) is not convex in the posterior.
The entry comparison turns on where posterior mass sits relative to the common bounds \(m(\pi),M(\pi)\).
